% ============================================================
\chapter{Objetivos concretos y metodología de trabajo}
\label{cap:objetivos}
% ============================================================

Este capítulo define los objetivos del trabajo y describe la metodología seguida
para alcanzarlos. En primer lugar se formula el objetivo general y los objetivos
específicos; los de naturaleza cuantificable incorporan criterios de éxito
medibles. A continuación se describen las seis fases del proceso de desarrollo
que estructuran el resto del trabajo.

% -----------------------------------------------------------
\section{Objetivo general}
% -----------------------------------------------------------

Diseñar, implementar y evaluar un sistema automático de predicción de cadenas de
ataque en infraestructuras \textit{Kubernetes} que complemente los enfoques de
análisis estático por manifiesto individual con el contexto estructural del
clúster, mediante la combinación de un clasificador de manifiestos individuales
(\textit{Random Forest}) con una red neuronal de grafos (\textit{Graph Attention
Network}) y optimización de ensemble mediante Algoritmo Genético. El sistema opera
principalmente sobre manifiestos \acs{YAML} estáticos, sin requerir acceso a un
clúster activo, si bien incorpora opcionalmente una vía de consulta directa al
clúster (OE6). Los objetivos específicos que concretan este objetivo general,
numerados OE1 a OE6, se detallan a continuación.

% -----------------------------------------------------------
\section{Objetivos específicos}
% -----------------------------------------------------------

\begin{description}
  \item[OE1.] Construir un conjunto de datos de manifiestos \textit{Kubernetes}
    etiquetados que combine múltiples fuentes públicas, contenga al menos
    2.000~manifiestos, y alcance una cobertura mínima del 2\,\% en la bandera de escape
    \texttt{TRUE\_HOST\_PID}, superando el 0,3\,\% de los repositorios reales no
    aumentados.

  \item[OE2.] Desarrollar una Capa~1 de clasificación de manifiestos individuales
    basada en \textit{Random Forest} con F1 macro superior a 0,85, que produzca
    \textit{risk scores} continuos en el rango [0,1] para usar como características
    de nodo en la Capa~2.

  \item[OE3.] Diseñar un esquema de grafo de clúster con tipos de arista
    semánticamente significativos para el dominio \textit{Kubernetes}:
    proximidad de directorio, alcance de privilegio, movimiento lateral mediante
    cuenta de servicio, co-espacio de nombres semántico y escalada \ac{RBAC}; su
    contribución se considera validada si la \ac{GAT} con tipos de arista iguala o
    supera a una \ac{GCN} sin ellos en F1 macro y Precisión@5 bajo idéntico
    protocolo de evaluación.

  \item[OE4.] Desarrollar una Capa~2 de clasificación de clústeres basada en \ac{GAT}
    que alcance Precisión@5 superior a 0,70 (media de validación cruzada de 5~pliegues)
    en la detección de cadenas de ataque, evaluada sobre el modelo desplegado con
    semilla fija documentada y verificada mediante análisis de robustez multi-semilla.

  \item[OE5.] Implementar una Capa~3 que optimice mediante Algoritmo Genético una
    función objetivo combinada de Precisión@5 y falsos positivos de clústeres
    limpios sobre las predicciones \textit{out-of-fold} de la validación cruzada,
    con criterios de éxito sobre el conjunto de test reservado:
    \ac{FPR}\_clean\,=\,0 y Precisión@5 no inferior a la del mejor componente
    individual del \textit{ensemble}.

  \item[OE6.] Empaquetar el sistema completo como una herramienta de línea de comandos
    (\textit{kubescan}) con soporte para análisis de directorios locales y,
    opcionalmente, consulta directa al \ac{API} de \textit{Kubernetes} mediante
    \textit{kubectl}; instalable en entornos Python 3.10+ mediante un único
    \texttt{pip install}.
\end{description}

% -----------------------------------------------------------
\section{Metodología del trabajo}
% -----------------------------------------------------------

El trabajo sigue un proceso iterativo de seis fases que combina el paradigma de
desarrollo de \textit{software} (requisitos, diseño, implementación, prueba) con el
proceso de ciencia de datos (adquisición, preparación, modelado, evaluación):

\textbf{Fase~1 — Adquisición de datos.} Descarga y normalización de manifiestos
\textit{Kubernetes} de cinco fuentes públicas complementarias (agrupadas en
cuatro categorías): el conjunto de datos
de Rahman et~al.\ \citep{Rahman2022} (repositorios reales etiquetados), BishopFox
BadPods \citep{BishopFox2021} (pods intencionadamente vulnerables), \textit{Kubernetes
Goat} \citep{KubernetesGoat} (escenarios de vulnerabilidad) y repositorios de
seguridad adicionales (CTFs, laboratorios y \textit{fixtures}). El criterio de
inclusión de repositorios adicionales fue la presencia de al menos una bandera de escape
activa (\texttt{HOSTPATH\_MOUNT}, \texttt{TRUE\_HOST\_PID}, etc.) en al menos un
manifiesto del directorio.

\textbf{Fase~2 — Extracción de características y construcción del grafo.} Extracción
de 28 características por manifiesto (25 indicadores binarios y 3 agregados
derivados) mediante análisis estático multi-herramienta
(Sección~\ref{sec:features_table}) y construcción del grafo de
clúster con cinco tipos de arista (Tabla~\ref{tab:edge_types}). En esta fase se
realizó también la validación de cobertura de la bandera \texttt{HOSTPATH\_MOUNT},
documentada en la Sección~\ref{sec:layer2}.

\textbf{Fase~3 — Entrenamiento de la Capa~1.} Entrenamiento de un clasificador
\textit{Random Forest} mediante validación cruzada de 5~pliegues para predecir la
presencia de misconfiguraciones individuales y generar \textit{risk scores}
calibrados, con los hiperparámetros fijos detallados en la
Sección~\ref{sec:capa1_diseno}.

\textbf{Fase~4 — Entrenamiento de la Capa~2.} Entrenamiento de una \ac{GAT} de 3~capas
con agrupamiento de grafos (\textit{graph pooling}) y validación cruzada estratificada
de 5~pliegues consciente de grupos y de familias de plantilla: las variantes
aumentadas siguen a su grafo base y las familias de plantilla con etiquetas mixtas
(p.~ej.\ \texttt{badpods\_*}) se asignan como unidades atómicas a una única
partición, nunca repartidas entre pliegues. La aumentación de datos se aplica sobre
los grafos de entrenamiento para mitigar la escasez de ejemplos de cadena de
ataque, y el \textit{checkpoint} de cada pliegue se selecciona por Precisión@5 de
validación.

\textbf{Fase~5 — Optimización de ensemble y evaluación.} Optimización de los pesos
del ensemble de tres componentes (\ac{RF}, \ac{GNN}, señal de escape) mediante
Algoritmo Genético sobre las predicciones \textit{out-of-fold} de la validación
cruzada de la Capa~2, aplicación de una restricción de viabilidad estructural en el
\textit{ranking} final (un grafo de un solo nodo no puede albergar una cadena de
ataque multi-salto), y evaluación final sobre el conjunto de test reservado desde
el inicio y no utilizado en ninguna fase de entrenamiento ni selección de
hiperparámetros.

\textbf{Fase~6 — Empaquetado y evaluación funcional y de usabilidad.} Empaquetado del
sistema como herramienta de línea de comandos instalable (\textit{kubescan}),
evaluación funcional de rendimiento e integración (latencia, formatos de salida,
resolución de \textit{checkpoints}) y estudio de usabilidad y aplicabilidad con
usuarios expertos mediante el cuestionario \ac{SUS} y un conjunto de tareas guiadas;
los protocolos y resultados de esta fase se detallan en el Capítulo~\ref{cap:evaluacion}.
